\documentclass{sysuredhead}

\RedHeadSetup{
  layout = normal,
  direction = down,
  authority = {示例市公共服务办公室},
  mark = {示例市公共服务办公室文件},
  document-number = {示公服发〔2026〕8号},
  title = {关于开展公共服务流程优化试点的通知},
  recipients = {各示例区公共服务部门，各有关单位：},
  issuer = {示例市公共服务办公室},
  date = {2026年9月2日},
  note = {联系人：示例甲，联系电话：0000-00000000。},
  copies = {示例市信息中心，示例市政务服务中心。},
  print-office = {示例市公共服务办公室综合处},
  print-date = {2026年9月2日},
  draft-label = {排版示例，不具公文效力},
  seal = placeholder
}

\begin{document}

\RedHeadMakeHead

为演示普通下行文的版式组件，现就虚构的公共服务流程优化试点事项通知如下。本文中的机关、编号、人员和事项均为虚构内容。

\section{明确试点范围}

各单位应结合示例场景梳理咨询、受理、办理和反馈环节，形成边界清楚、便于检查的流程清单。

\subsection{统一服务事项}

同一事项的名称、条件、材料和办理时限应保持一致。确需调整的，应当同步更新面向服务对象的说明。

\subsubsection{完善反馈渠道}

对试点中发现的问题分类记录，并明确负责人员和处理期限。

\paragraph{保护测试数据}

演示数据不得包含个人敏感信息、工作秘密或者其他受限信息。

\newpage

\section{组织实施}

各单位应在虚构试点周期内完成自查、联调和复盘。涉及跨部门协作的事项，由示例牵头单位汇总协调。

\section{工作要求}

本示例只展示排版结构，不构成任何现实工作部署。正式公文应由有权机关依照程序审核、签发和用印。

\begin{RedHeadAttachments}
  \item 公共服务流程检查表
  \item 试点问题记录表
\end{RedHeadAttachments}

\RedHeadMakeClosing

\begin{RedHeadAttachment}{1}{公共服务流程检查表}
\section*{一、基本信息}

填写事项名称、负责单位和检查日期。

\section*{二、检查内容}

依次记录咨询、受理、办理和反馈环节的示例结果。
\end{RedHeadAttachment}

\begin{RedHeadAttachment}{2}{试点问题记录表}
本页用于演示第二个附件及连续页码。请勿填写真实敏感信息。
\end{RedHeadAttachment}

\RedHeadMakeImprint

\end{document}
